\documentclass[11pt]{article}
\usepackage[margin=1in]{geometry}
\usepackage{amsmath,amssymb,amsthm}
\usepackage{hyperref}
\hypersetup{hidelinks}

\newtheorem{theorem}{Theorem}
\newtheorem{lemma}[theorem]{Lemma}
\newtheorem{corollary}[theorem]{Corollary}
\theoremstyle{definition}
\newtheorem{definition}{Definition}
\theoremstyle{remark}
\newtheorem{remark}{Remark}

\setlength{\parskip}{4pt plus 1pt minus 1pt}
\setlength{\parindent}{0pt}

\title{Evaluation-Driven Descent, Encoding Permanence,\\and the Structural Invariants of Intelligence}
\author{Patrick D.\ McCarthy}
\date{}

\begin{document}
\maketitle

\begin{abstract}
We establish four linked results for any entity with structure
$(K, A, \sigma, \gamma)$ persisting under bounded active context,
survival pressure, and positive novelty rate $\nu(E) > 0$.
First, agency ($\gamma > 0$) necessarily produces evaluation-driven
state revision---monotone local descent on expected uncertainty
$H(W \mid K)$---where each informative update reduces $H(W \mid K)$
by at most the information gain $G(a, K)$, with equality under exact
Bayesian conditioning. Under non-stationary task distribution,
descent is local and instantaneous rather than globally convergent.
Second, the certainty threshold for encoding an action at level $L_i$
is $\theta_i = 1 - \varepsilon / C_i$, which derives the learning
rate hierarchy from cost-of-failure differences rather than
assumption.
Third, separating the action space $A$ from the higher-order function
space $M$ is selectively favored under bounded context: collapsing them
is possible but increasingly costly under positive novelty rate, so
competition eliminates collapsed architectures.
Fourth, $\gamma > 0$ and learning efficiency $\eta_M > 0$ are the
unique structural invariants of persistence: the free-context
exhaustion argument shows that $\gamma = 0$ entities are eliminated
under positive novelty rate; without $\eta_M > 0$, active context
fills without relief and inference degrades to failure. Every other
property can equal zero for some persistent entity. Therefore
$I(E) = \gamma \cdot \eta_M$---the rate of adaptation to novelty---is
the unique measure of intelligence, derived not stipulated.
\end{abstract}

% ------------------------------------------------------------------
\section{Introduction}
\label{sec:intro}
% ------------------------------------------------------------------

This paper completes a chain of results deriving the necessary
architecture of intelligence from a single generative constraint:
bounded active context.

Paper~1 \cite{McCarthy2026a} established that any entity accumulating
propositions under bounded context necessarily maintains a directed
graph with typed edges---where ``necessary'' means selected for under
competition, not logically impossible to avoid (see \cite{McCarthy2026a},
Remark~1). Paper~2 \cite{McCarthy2026b} established that
bounded context under growing action count requires convergence toward
escalation. Paper~3 \cite{McCarthy2026c} established that bounded context forces the
normalization of shared propositions from action into a separated $K$,
with an asymmetric retention criterion.

Three questions remain. First: what dynamics govern the updating
process when an entity with agency acts on the world and receives
feedback? Second: what determines the threshold at which a proven
action earns promotion to a more permanent encoding level? Third:
what properties must every persisting entity have?

We answer all three. Theorem~\ref{thm:descent} establishes that
agency ($\gamma > 0$) with feedback produces evaluation-driven state
revision---monotone local descent on expected uncertainty
$H(W \mid K)$. Definition~\ref{def:encoding} establishes that the
certainty threshold for encoding at level $L_i$ is
$\theta_i = 1 - \varepsilon / C_i$, deriving the learning rate
hierarchy from cost-of-failure differences.
Theorem~\ref{thm:type} establishes that separating the action space
$A$ from the higher-order function space $M$ is selectively favored:
collapsing them is increasingly costly under positive novelty rate.
Theorem~\ref{thm:invariants} establishes that $\gamma > 0$ and
learning efficiency $\eta_M > 0$ are the unique structural invariants
of persistence, and Corollary~\ref{cor:measure} derives
$I(E) = \gamma \cdot \eta_M$ as the unique intelligence measure.

Several lines of prior work converge on these questions.
Dupoux, LeCun, and Malik \cite{DupouxLeCunMalik2026} propose a
bilevel optimization framework but treat the learning rate
relationship between levels as a design parameter rather than deriving
it, and deliberately avoid defining intelligence. Hutter
\cite{Hutter2005} and Legg and Hutter \cite{LeggHutter2007} define
intelligence as average performance over computable environments;
we define it as rate of improvement, independent of current
performance level. Chollet \cite{Chollet2019} measures
skill-acquisition efficiency, structurally close to our $\eta_M$ but
treated as a benchmark metric rather than a structural invariant.
Schmidhuber \cite{Schmidhuber2010} identifies curiosity as compression
progress but treats it as a design objective rather than a necessary
condition of persistence. Friston's free energy principle
\cite{Friston2010} derives minimization behavior from
self-organization but does not distinguish $\gamma = 0$ from
$\gamma > 0$ within its framework. Cover and Thomas \cite{Cover2006}
provide the information-theoretic foundations that anchor our
convergence argument. The contribution of this paper is to derive,
via selection and exhaustion arguments, the unique structural
invariants of persistence and the intelligence measure they determine.

% ------------------------------------------------------------------
\section{Formal Framework}
\label{sec:framework}
% ------------------------------------------------------------------

We introduce the definitions required for the four main results,
extending the framework of \cite{McCarthy2026a}.

\begin{definition}[Entity, World, Uncertainty]
\label{def:entity}
An entity $E = (K, A, \sigma, \gamma)$ consists of a knowledge base $K$
(propositions about world $W$), an action space $A$ (executable
actions), a selection mechanism $\sigma$ that chooses actions given
context (which presupposes sensing capacity; see
Theorem~\ref{thm:invariants}), and an agency parameter $\gamma$. The world state space $W$ is a
measurable space; the true state $w \in W$ is not directly observable.
The \emph{uncertainty} of $E$ about state $w$ given knowledge $K$ is:
\[
U(w, K) = -\log P(w \mid K)
\]
The expected uncertainty is the conditional entropy:
\[
H(W \mid K) = \mathbb{E}_w[U(w, K)] = -\sum_w P(w \mid K) \log P(w \mid K)
\]
\end{definition}

\begin{definition}[Agency]
\label{def:agency}
Agency $\gamma \in [0, 1]$ is the probability that the entity engages its
adaptive capacity $M$ in response to a learning signal. Equivalently,
$\gamma$ is the response gate: the fraction of environmental signals that
trigger an update to $(A, K)$.

The operational signature of $\gamma > 0$ is that the entity bounds
uncertainty beyond what immediate survival requires, freeing context
capacity for novel states. An entity with $\gamma = 0$ responds to no
signal and accumulates no free context. An entity with $\gamma > 0$
engages $M$ on signals, resolves uncertainty, promotes actions to lower
encoding levels, and thereby maintains $C_{\mathrm{free}} > 0$
(Definition~\ref{def:free}).

The distinction between $\gamma = 0$ and $\gamma > 0$ is not between
entities that face threats and those that do not. It is between entities
that build reserve capacity for states they have not yet encountered and
those that do not.
\end{definition}

\begin{definition}[Action Feedback]
\label{def:feedback}
For action $a$ executed in world state $w$, the \emph{feedback}
$\varphi(a, w)$ is the observable consequence. The entity updates its
knowledge base: $K_a = K \cup \{p_a\}$, where $p_a$ is the
proposition encoding the feedback. The \emph{prediction-error signal} is:
\[
\delta = \varphi(a, w) - \mathbb{E}[\varphi(a, \cdot) \mid K]
\]
\end{definition}

\begin{definition}[Information Gain]
\label{def:gain}
When entity $E$ executes action $a$ in world state $w$, it receives
feedback $\varphi(a, w)$ encoded as proposition $p_a$
(Definition~\ref{def:feedback}). The \emph{information gain} is the
mutual information between the feedback proposition and the world
state, conditioned on current knowledge:
\[
G(a, K) = I(p_a; W \mid K) \geq 0
\]
Non-negativity follows from the definition of mutual information
\cite{Cover2006}. $G = 0$ if and only if $p_a$ provides no
information about $W$ beyond what $K$ already contains.
\end{definition}

\begin{definition}[Novelty Rate]
\label{def:novelty}
The \emph{novelty rate} $\nu(E)$ of entity $E$ in world $W$ is the rate
at which $E$ encounters world states $w$ where $U(w, K) > 0$:
\[
\nu(E) = \lim_{T \to \infty} \frac{1}{T}
\bigl|\{t \leq T : U(w_t, K(t)) > 0\}\bigr|
\]
Novelty rate is a property of the gap between $W$ and $K$, not of $W$
alone. A world with high entropy but an entity whose $K$ covers all
reachable states has $\nu = 0$. The distinction matters: entropy is
a property of the environment; novelty is a property of the
entity--environment coupling. The elimination argument
(Theorem~\ref{thm:invariants}) requires that novelty include
\emph{survival-relevant} states: states where $U(w, K)$ can exceed
$U_{\mathrm{lethal}}$ unless $K$ is updated. Under positive entropy in
$W$, the world state distribution drifts, and previously non-lethal gaps
in $K$ eventually become lethal. The assumption is that $\nu > 0$
includes eventual survival-relevant novelty, not merely benign
uncertainty.
\end{definition}

\begin{definition}[Encoding Hierarchy]
\label{def:hierarchy}
The encoding space $\mathcal{E}(c, r)$ is a continuous space
parameterized by cost $c$ and reversibility $r$. Within this space,
interfaces $L_0, L_1, \ldots, L_n$ define discrete encoding levels
with increasing cost and reversibility:
\begin{itemize}
\item $L_0$ (genomic): highest cost of error, lowest reversibility,
   slowest update rate.
\item $L_n$ (active context): lowest cost of error, highest
   reversibility, fastest update rate.
\end{itemize}
Escalation \cite{McCarthy2026b} pushes proven actions
upward toward more permanent levels. The present paper derives the
threshold that governs this promotion.
\end{definition}

\begin{definition}[Free Context]
\label{def:free}
The \emph{free context} of entity $E$ at time $t$ is the active context
capacity not consumed by unresolved content:
\[
C_{\mathrm{free}}(t) = C_n - |\mathrm{ctx}_{\mathrm{unresolved}}(t)|
\]
where $\mathrm{ctx}_{\mathrm{unresolved}}(t)$ is the set of
propositions in active context that have not yet been encoded at a lower
level or resolved by $M$. Free context is the entity's reserve capacity
for handling novel states. The architecture derived in
Papers~1--3---graph structure, escalation, retention criterion---exists
to maximize $C_{\mathrm{free}}$: each result describes a mechanism that
moves resolved content out of $L_n$, preserving capacity for what the
entity has not yet encountered.
\end{definition}

\begin{definition}[Higher-Order Function Space $M$]
\label{def:M}
Let $A : K \times W_{\mathrm{obs}} \to \Delta \mathcal{A}$ be the
action space of informed actions, where $W_{\mathrm{obs}}$ denotes
observable world states and $\Delta \mathcal{A}$ the simplex over
available actions. The \emph{higher-order function space} is
\[
M : (A, K) \times L \;\to\; (A, K)
\]
where $L$ is a learning signal (feedback from $W$ following action).
$M$ is categorically different from $A$: elements of $A$ map knowledge
and observations to action distributions; elements of $M$ map
action--knowledge pairs and learning signals to updated
action--knowledge pairs.

$M$ contains whatever adaptive capacity has not yet been committed to
specific knowledge or action. As the entity learns, components of $M$
resolve into concrete elements of $K$ or $A$---the same normalization
by which shared propositions precipitate from action into knowledge
\cite{McCarthy2026c}. What remains in $M$ is the entity's capacity
to produce further such resolutions.

The hierarchy is open upward: $M^{(2)}$ operates on $M^{(1)}$, yielding
meta-learning; $M^{(n+1)}$ operates on $M^{(n)}$. No natural
termination point is established (see Section~\ref{sec:open}).
\end{definition}

\begin{definition}[$M$ Efficiency]
\label{def:eta}
Given a learning event where $(A', K') = M((A, K), L)$, define the
\emph{learning efficiency}:
\[
\eta_M \;=\; \mathbb{E}_{L \mid \text{engaged}}\!\left[
\frac{\max_{a' \in A'} G(a', K') - \max_{a \in A} G(a, K)}{|L|}
\right]
\]
The expectation is taken over learning signals $L$ \emph{conditional on
the entity engaging $M$} (i.e., conditional on the gate $\gamma$
opening). This conditioning is essential: $\gamma$ may be directed,
so $\mathbb{E}[{\cdot} \mid \text{engaged}]$ may differ from the
unconditional expectation. Conditioning on engagement ensures that the
product $\gamma \cdot \eta_M$ is the expected improvement rate by the
law of total expectation, regardless of whether $\gamma$ is
content-blind or directed.

Both numerator and denominator are measured in bits; $\eta_M$ is
dimensionless. The three regimes:
\begin{itemize}
\item $\eta_M > 0$: the entity learns---grounding improves per unit
      of learning signal, conditional on engagement.
\item $\eta_M = 0$: engagement produces no improvement.
\item $\eta_M < 0$ temporarily: exploration, creative destruction,
      locally irrational steps that may increase long-run survival.
      Persistently $\eta_M < 0$: the entity does not persist.
\end{itemize}
\end{definition}

\begin{remark}[Measurement]
\label{rem:measurement}
Computing $\eta_M$ requires evaluating $G(a, K) = I(p_a; W \mid K)$,
which depends on the true joint distribution $P(p_a, W, K)$. The entity
cannot compute its own $\eta_M$; it would need the very information
it is trying to learn. Both $\gamma$ and $\eta_M$ are measurable by an
external observer with access to the entity's trajectory and outcomes.
\end{remark}

\begin{definition}[Intelligence]
\label{def:intelligence}
For an entity $E$ with agency $\gamma$ and learning efficiency $\eta_M$:
\[
I(E) = \gamma \cdot \eta_M
\]
Intelligence is the rate of adaptation to novelty: the expected rate at
which the entity reduces uncertainty about novel states. This is stated
here as a definition. Theorem~\ref{thm:invariants} establishes that
$\gamma > 0$ and $\eta_M > 0$ are the unique structural invariants
of persistence under positive novelty rate.
\end{definition}

\begin{definition}[Persistence]
\label{def:persistence}
An entity $E$ \emph{persists} if it survives for all $t$: there is no
finite $T$ at which $U(w_T, K(T)) > U_{\mathrm{lethal}}^d$ terminates
the entity. Persistence is indefinite survival, not survival to some
fixed horizon.
\end{definition}

% ------------------------------------------------------------------
\section{Monotone Descent on Uncertainty}
\label{sec:descent}
% ------------------------------------------------------------------

\begin{theorem}[Evaluation-Driven State Revision]
\label{thm:descent}
Let $E = (K, A, \sigma, \gamma)$ be an entity with $\gamma > 0$ persisting in
world $W$ under bounded active context $C_n$.

\textbf{Part I.} Every executed action generates a prediction-error
signal, and $\gamma > 0$ guarantees engagement with that signal,
producing evaluation-driven state revision---monotone local descent on
expected uncertainty $H(W \mid K)$ within the entity's model. Under
competition (Remark~1 of \cite{McCarthy2026a}), entities that fail to
descend are eliminated: descent is selectively necessary, not merely
optimal.

\textbf{Part II.} Under positive novelty rate $\nu(E) > 0$
(Definition~\ref{def:novelty}), descent never terminates globally:
the entity continues to encounter states where $U(w,K) > 0$, so
actions with positive information gain always exist.
\end{theorem}

\begin{proof}
\textbf{Part I: Forced Signal, Response Gate, and Convergence.}

Every executed action $a$ generates feedback $\varphi(a, w)$ from the
world. The prediction error $\delta = \varphi(a, w) -
\mathbb{E}[\varphi(a, \cdot) \mid K]$ is produced at every action
regardless of $\gamma$.

Agency $\gamma$ is the response gate. An entity with $\gamma = 0$ receives
every signal $\delta$ but never updates: $(A, K)$ remain static. An
entity with $\gamma > 0$ engages $M$ with probability $\gamma$ on each
signal, integrating the feedback into $K$ by forming
$K' = K \cup \{p_a\}$. $K$ is a set of propositions, not a parameter
vector, so $\nabla_K U$ is undefined: there is no differentiable loss
surface over $K$. The descent property is instead established directly
via the Lyapunov function $H(W \mid K)$.

\emph{Step 1: Each update reduces the entity's estimate of $H(W \mid K)$.}
When the entity encounters a novel state and executes action $a$, the
feedback $p_a$ is informative: $G(a, K) = I(p_a; W \mid K) > 0$.
When $M$ integrates this feedback, forming $K' = K \cup \{p_a\}$:
\[
H(W \mid K') = H(W \mid K) - G(a, K)
\]
This identity follows from the chain rule for mutual information under
exact Bayesian conditioning: $H(W \mid K') = H(W \mid K, p_a)$.

\textbf{The distinction between model and reality.}
The descent is on the entity's \emph{model} of $W$, not on the true
state of $W$. Adding a proposition to $K$ refines the posterior within
the entity's model---this is monotone by construction, since adding a
constraint to the joint likelihood $L_K$ cannot broaden the posterior
(see the model-vs-reality remark in \cite{McCarthy2026c}, Lemma~1).
But the proposition may be wrong: the evidence may be noise, the
likelihood function may be mis-specified, or the world may have shifted
since the evidence was collected. In such cases, the entity's model
diverges from reality even as its \emph{estimated} $H$ decreases.

This is not a gap in the argument but the precise content of the
persistence conditions established in Theorem~\ref{thm:invariants}.
The descent on the entity's model is guaranteed at every step. Whether
that descent tracks reality depends on the quality of evidence
integration---which is $\eta_M$. An entity with
$\mathbb{E}_t[\eta_M] > 0$ improves its model's correspondence with
$W$ on average over its trajectory. An entity whose model
systematically diverges ($\mathbb{E}_t[\eta_M] < 0$ persistently) is
exactly the failure case that Theorem~\ref{thm:invariants}, Part~2
eliminates. Temporary divergence ($\eta_M < 0$ at individual steps)
is compatible with persistence and is the mechanism of exploration
(Remark~\ref{rem:exploration}).

The Lyapunov argument therefore has two levels: (i)~within the model,
descent is monotone at every informative step; (ii)~across the
model-reality interface, descent tracks reality only in expectation,
guaranteed by $\mathbb{E}_t[\eta_M] > 0$. Level~(i) is proved here.
Level~(ii) is the persistence condition.

\emph{Step 2: Monotone convergence within the model.} Consider the
sequence $\{H_t\} = \{H(W \mid K(t))\}$ indexed by informative update
steps, where $H$ is the entropy \emph{as estimated by the entity's
model}. By Step~1, $\{H_t\}$ is strictly decreasing at each
informative step. By definition, $H(W \mid K) \geq 0$, so the sequence
is bounded below. By the monotone sequence theorem, $\{H_t\}$
converges. Since the sequence is non-increasing and convergent:
\[
\sum_{t=0}^{\infty} G(a_t, K(t)) \leq H(W \mid K(0)) < \infty
\]
Therefore $G(a_t, K(t)) \to 0$ as $t \to \infty$. The entity
converges to a state where available actions yield diminishing
informational returns \emph{within its current model}.

\textbf{Part II: Non-Termination Under Positive Novelty.}

When $\nu(E) > 0$, the entity encounters states $w$ with
$U(w, K(t)) > 0$ at positive rate. Each such encounter produces an
action $a$ with $G(a, K(t)) > 0$. Therefore at every finite $t$ there
exists some $a'$ with $G(a', K(t)) > 0$. Descent converges locally
(Part~I) but never terminates globally.

Non-termination does not require unbounded world complexity. A finite
world with $\nu(E) > 0$ suffices: even if $W$ is finite, the
combinatorial space of entity-environment interactions exceeds what $K$
can index under bounded $C_n$.

\emph{Interaction with the curation regime.} The convergence argument
in Part~I assumes $K$ grows monotonically. In
the curation regime established in \cite{McCarthy2026c}, $|K|$ is
bounded and every addition requires displacement. The
Lyapunov function must be refined. Raw $H(W \mid K)$ is the wrong
quantity: displacing a proposition about an irrelevant domain in favor
of one about a survival-critical domain could increase $H(W \mid K)$
overall while improving the entity's survival prospects. The correct
Lyapunov function is \emph{survival-weighted expected uncertainty}:
\[
H_{\mathcal{T}}(W \mid K) = \mathbb{E}_{w \sim \mathcal{T}}[U(w, K)]
\]
where $\mathcal{T}$ is the entity's task distribution over
survival-relevant world states.

Under the cost-benefit displacement criterion of \cite{McCarthy2026c},
proposition $p$ is displaced by $q$ only when
\[
  \mathbb{E}[B(q, \mathcal{T})] - c_{\text{retain}}(q, K)
  \;>\;
  \mathbb{E}[B(p, \mathcal{T})] - c_{\text{retain}}(p, K),
\]
so each displacement reduces $H_{\mathcal{T}}$.

$\mathcal{T}$ itself is non-stationary: as the entity learns and the
environment shifts, the distribution of survival-relevant states
changes. Under non-stationary $\mathcal{T}$, the global Lyapunov
convergence argument does not apply: $H_{\mathcal{T}_t}$ at step $t$
and $H_{\mathcal{T}_{t+1}}$ at step $t+1$ evaluate different functions
when $\mathcal{T}$ drifts. What the curation regime establishes is
\emph{local, instantaneous} descent: each displacement reduces
$H_{\mathcal{T}}$ evaluated against the current $\mathcal{T}$ at the
moment of displacement. This is a weaker claim than global
convergence, but it is the operationally relevant one: the entity
cannot plan against future task distributions it has not yet
encountered. The non-stationarity of $\mathcal{T}$ is precisely what
drives non-termination (Part~II): the target keeps moving, so the
entity never reaches $H_{\mathcal{T}} = 0$ globally.

\textbf{Remark.} An entity with $\gamma = 0$ never updates. As $W$
changes, its fixed $K$ becomes increasingly inaccurate: $U(w_t, K)$
grows until it exceeds $U_{\mathrm{lethal}}$. Selection removes such
entities. Therefore $\gamma > 0$ is necessary for persistence, a claim
established fully in Theorem~\ref{thm:invariants}.
\end{proof}

% ------------------------------------------------------------------
\section{Encoding Permanence}
\label{sec:encoding}
% ------------------------------------------------------------------

\begin{definition}[Encoding Permanence Scales with Rigor]
\label{def:encoding}
Let $L_0, \ldots, L_n$ be encoding levels with associated failure
costs $C_0 > C_1 > \cdots > C_n$ and reversibility costs
$r_0 > r_1 > \cdots > r_n$. Let $\varepsilon$ be the maximum
acceptable expected error cost. Then the minimum certainty threshold
for encoding an action at level $L_i$ is:
\[
\theta_i = 1 - \frac{\varepsilon}{C_i}
\]
\end{definition}

\textbf{Derivation.}
If $\theta$ is the certainty that action $a$ is correct, the expected
error cost is $(1 - \theta) \cdot C_i$. Encoding is safe when:
\[
(1 - \theta) \cdot C_i \leq \varepsilon
\quad\Longrightarrow\quad
\theta \geq 1 - \frac{\varepsilon}{C_i}
\]
The minimum safe threshold is $\theta_i = 1 - \varepsilon / C_i$.
This is arithmetic, not a deep result. The substantive content is the
\emph{identification}: that the certainty threshold is determined by
cost-of-failure, not by a free parameter. Given that identification,
the formula follows in one line.

\textbf{Consequence.} At $L_0$ (genomic), $C_0$ is the cost of lineage
extinction, so $\theta_0 \approx 1$: near-certainty is required.
At $L_n$ (active context), $C_n$ is small, so $\theta_n$ is
meaningfully less than $1$: the expected benefit of trying exceeds the
expected cost of failure at moderate confidence.

\textbf{Error propagation through the hierarchy.} Under escalation
\cite{McCarthy2026b}, errors partition into \emph{handled errors},
resolved at level $L_i$, and \emph{unhandled exceptions}, which
escalate. The error rate seen by $L_{i+1}$ is the rate of genuinely
irreducible errors, not a compounded Bernoulli product of lower-level
rates. The thresholds $\theta_i$ are independent by construction.

\subsection{Deriving the Learning Rate Hierarchy}

The learning rate at level $L_i$ is inversely related to the rigor
threshold. Each trial of action $a$ yields a binary outcome (success or
failure), so the empirical success rate $\hat{p}$ after $n$ trials is
the sample mean of $n$ independent Bernoulli random variables in
$[0,1]$. Hoeffding's inequality \cite{Hoeffding1963} applies directly
to sample means of bounded random variables:
\[
P(|\hat{p} - p| \geq \delta) \leq 2\exp(-2n\delta^2)
\]
The entity requires confidence $\theta_i$ that $\hat{p}$ is within
$\delta$ of $p$. Setting the failure probability
$2\exp(-2n\delta^2) \leq 1 - \theta_i = \varepsilon / C_i$:
\[
n_i \;\geq\; \frac{\log(2 C_i / \varepsilon)}{2\delta^2}
\]
The effective learning rate is inversely proportional to $n_i$:
\[
\eta_i \;\propto\; \frac{1}{n_i}
\;\propto\; \frac{\delta^2}{\log(C_i / \varepsilon)}
\]

\begin{remark}[Level-independence of $\delta$]
\label{rem:delta}
The deviation tolerance $\delta$ measures how close the empirical
success rate must be to the true rate---a property of measurement
accuracy, not of the stakes. Assuming $\delta$ constant across levels is
the conservative choice: a rational entity might demand tighter
precision at $L_0$ ($\delta_0 < \delta_n$), which would make
$\eta_{\mathrm{evo}}$ even smaller relative to $\eta_{\mathrm{ind}}$,
strictly strengthening the hierarchy. Level-independent $\delta$
produces the \emph{weakest} version of the learning rate separation.
\end{remark}

Since $\theta_0 \gg \theta_n$, it follows that
$\eta_{\mathrm{evo}} \ll \eta_{\mathrm{ind}}$. The learning rate
difference between evolution and individual learning is \emph{derived}
from the cost-of-failure structure, not assumed as a design parameter.

\begin{definition}[Certainty]
\label{def:certain}
Action $a$ is \emph{certain} at level $L_i$ when
$P(\text{$a$ is adaptive} \mid K) \geq \theta_i$ across
sufficient variation in $w$. Once $\theta_i$ is met, the action is
encoded at $L_i$ and no longer consumes $L_n$ resources, freeing
active context for the next epistemic frontier.
\end{definition}

\begin{corollary}[Bounded Adaptation]
\label{cor:bounded}
Encoding at $L_i$ is durable, not literally irreversible:
environmental shift can render it maladaptive, but correction requires
effort proportional to $r_i$. Individual learning at $L_n$ discovers
actions via rapid evaluation-driven descent ($\eta_{\mathrm{ind}}$
large). Encoding permanence determines which promote to $L_0$: only
those meeting $\theta_0 \approx 1$. This is the Baldwin effect
\cite{Hinton1987}, here derived rather than observed.
\end{corollary}

% ------------------------------------------------------------------
\section{Convergence to $A$/$M$ Stratification}
\label{sec:type}
% ------------------------------------------------------------------

\begin{theorem}[Convergence to $A$/$M$ Stratification]
\label{thm:type}
Systems that collapse the action space $A$ and the learning space $M$
into a single parameter space pay an increasing cost under bounded
context. Under bounded $C_n$ with positive novelty rate, stratified
$A$/$M$ separation is selectively favored: collapsed systems pay
higher costs and are eliminated under competition
(\cite{McCarthy2026a}, Remark~1).
\end{theorem}

\begin{proof}
The type signatures of $A$ and $M$ are distinct:
\begin{align*}
A &: K \times W_{\mathrm{obs}} \to \Delta \mathcal{A} \\
M &: (A, K) \times L \to (A, K)
\end{align*}
$A$ maps knowledge and observations to action distributions. $M$ maps
action--knowledge pairs and learning signals to updated
action--knowledge pairs. These are different functions serving
different roles. The question is not whether a system \emph{must}
separate them---it need not---but what happens when it does not.

\textbf{The degenerate case.}
A system with no separation at all has a single function type
$T : W_{\mathrm{obs}} \to \Delta\mathcal{A}$.
This is the case $K = \varnothing$: no propositions are separated
from the action policies, so every action node carries its own copy
of whatever world-model information it needs. Shared propositions
are duplicated across nodes rather than extracted---the unfactored
state of \cite{McCarthy2026a}. The system functions but its nodes
are heavy: as the action repertoire grows, the redundant weight
grows with it.

Factoring $K$ out of $A$ converts that redundant node weight into
edges---each action links to shared propositions rather than
storing its own copy. Separating $M$ does the same for learning:
update logic is extracted once rather than embedded in every action.
This is exactly the normalization argument of \cite{McCarthy2026c}
applied to function types rather than propositions.

More generally, $M$, $A$, and $K$ can all be unified into a single
parameter space. The question is efficiency under bounded context,
not logical possibility.

\textbf{The cost of collapsing $A$ and $M$.}
A system that uses the same parameters for both action selection and
learning faces a structural conflict: updating parameters to
incorporate new learning (the $M$ function) modifies the same
substrate that encodes stable action policies (the $A$ function).
Without stratification, every learning update is an unscoped
modification to the action space. This is precisely the mechanism
identified in \cite{McCarthy2026a}: without stored adjacency between
what is being updated and what depends on it, updates cascade
uncontrollably.

The empirical manifestation is catastrophic forgetting
\cite{Kirkpatrick2017}: neural networks that learn new tasks degrade
performance on previously learned tasks because the weight updates
that encode new knowledge overwrite the weights that encode old
action policies. The network cannot scope the update because it lacks
the structural separation between ``what I know how to do'' ($A$)
and ``how I change what I know'' ($M$).

\textbf{Stratification under bounded context.}
Under bounded $C_n$, the cost of catastrophic forgetting is not merely
wasteful but fatal. Each unscoped update that degrades a proven action
forces that action back to $L_n$ for re-learning, consuming the
context capacity that the entity needs for novel problems. Under
positive novelty rate $\nu > 0$, novel states arrive continuously.
An entity that must re-learn previously encoded actions while
simultaneously handling novel states exhausts $C_{\mathrm{free}}$.

Stratification resolves this: $A$ encodes stable action policies at
lower encoding levels (protected from learning updates), while $M$
operates at $L_n$ on the current learning frontier. Updates to $M$
do not modify $A$; modifications to $A$ occur only through the
promotion mechanism of Definition~\ref{def:encoding}, gated by the
certainty threshold $\theta_i$. This is the same pattern as the
series: not a logical impossibility of collapsing $A$ and $M$, but
convergence under bounded-context pressure toward separation because
the collapsed alternative becomes increasingly costly.

\textbf{Empirical convergence.}
Systems that approximate $A$/$M$ stratification---elastic weight
consolidation \cite{Kirkpatrick2017}, progressive neural networks,
complementary learning systems \cite{McClelland1995}---outperform
flat architectures on continual learning precisely because they
protect encoded actions from learning interference. The trajectory
from flat parameter spaces toward stratified architectures is
the convergence this theorem predicts.
\end{proof}

\begin{corollary}[Distinct Roles]
\label{cor:roles}
The selectively favored separation establishes distinct functional roles:
$A$ navigates $W$ given $K$; $M$ evolves $(A, K)$ in response to
learning signals. Concrete instances: neuroplasticity ($M$ on $A$),
insight ($M$ on $K$), habit formation ($M$ promoting from $L_n$ to
lower levels per Definition~\ref{def:encoding}), evolution ($M$ on $M$
at the most stringent rigor threshold).
\end{corollary}

% ------------------------------------------------------------------
\section{Necessity of $\gamma$ and $\eta_M$}
\label{sec:necessity}
% ------------------------------------------------------------------

\begin{theorem}[Structural Invariants of Persistence]
\label{thm:invariants}
For any entity $E$ persisting (Definition~\ref{def:persistence}) under
bounded active context $C_n$, survival threshold
$U_{\mathrm{lethal}}$, with positive novelty rate $\nu(E) > 0$:
\begin{enumerate}
\item[(i)] $\gamma(E) > 0$ necessarily.
\item[(ii)] $\eta_M(E) > 0$ necessarily (in expectation over the
      entity's trajectory).
\item[(iii)] Every other scalar property of $E$ except positive
      sensing capacity (a precondition of $\sigma$) can equal zero
      for some persistent entity.
\item[(iv)] Therefore $I(E) = \gamma \cdot \eta_M$ is the natural measure
      of intelligence: the rate of adaptation to novelty, identified
      by what positive $\nu(E)$ selects for.
\end{enumerate}

When $\nu(E) = 0$, the entity faces no novel states and $\gamma = 0$ is
viable: no free context is needed because nothing unexpected arrives.
Such entities persist without intelligence. Intelligence is what positive
novelty rate selects for.
\end{theorem}

\textbf{Scope note.} This theorem applies to entities with the
structure $E = (K, A, \sigma, \gamma)$. This is a presupposition, not a
conclusion: the theorem does not derive the \emph{existence} of agency
from non-agentive precursors. What it establishes is an
\emph{elimination result}: among entities that have this structure,
those with $\gamma = 0$ do not persist under $\nu > 0$. The
contribution is the \emph{mechanism}---free-context exhaustion---not
the origin of $\gamma > 0$ (which remains open;
Section~\ref{sec:open}). Inert objects lack this structure entirely
and are outside the theorem's domain.

\begin{proof}

\textbf{Part 1: Necessity of $\gamma > 0$: World Drift and Context
Exhaustion.}

Suppose $\gamma(E) = 0$. By Definition~\ref{def:agency}, the entity never
engages $M$: $(A(t), K(t)) = (A(0), K(0))$ for all $t$.

The entity has $\nu(E) > 0$: it encounters states where $U(w, K) > 0$
at positive rate. Under positive entropy, the world drifts: states that
$K(0)$ was calibrated for become less frequent, and states not covered
by $K(0)$ become more frequent. The entity must still act to survive.
Each encounter with a novel state---one where no lower encoding level
has an adequate response---requires real-time processing at $L_n$,
consuming active context capacity.

An entity with $\gamma > 0$ resolves these encounters: it updates $K$,
promotes proven actions to lower levels, and frees $L_n$ for the next
problem. Each resolution converts a novel state into a state handled
below $L_n$. An entity with $\gamma = 0$ cannot resolve anything.
The same types of novel states keep arriving and keep requiring $L_n$
processing. As the world drifts further from $K(0)$, an increasing
fraction of the entity's encounters are novel---$K(0)$ covers less of
current $W$---so the concurrent demand on $L_n$ grows.
$C_{\mathrm{free}}(t) \to 0$ not because $K$ grows (it is fixed) but
because the fraction of encounters requiring active context increases
while no encounter is ever resolved.

When $C_{\mathrm{free}} = 0$, the entity cannot load the relevant
subgraph $K^{[a]}$ for any new problem: inference quality
$\rho \to 0$. The entity encounters $w$ where
$U(w, K(0)) > U_{\mathrm{lethal}}^d$ and does not survive.

The argument is self-contained at the individual level and requires no
competition. For any initial $K(0)$, because $\nu > 0$ and the world
drifts, there exists a finite time $T^*$ at which the fraction of
encounters requiring $L_n$ exceeds what $C_n$ can service. After $T^*$,
the entity cannot process novel states at the rate they arrive. Since
$\nu > 0$ continues, the entity eventually encounters a lethal state.
No competitor is required: the environment alone is sufficient.

The critical distinction is between $\nu = 0$ and $\nu > 0$. An entity
in a zero-novelty environment (a thermostat in a stable room) faces no
context pressure and persists with $\gamma = 0$. The framework does not
predict it dies. It predicts it does not need intelligence.

Positive entropy in $W$ is the upstream cause of positive novelty rate:
energy gradients produce states not contained in any finite $K$
\cite{Cover2006, Ashby1956}. But entropy is a property of the
environment; novelty rate is a property of the entity--environment
coupling. Entropy guarantees that $\nu > 0$ eventually for any entity
whose niche is not permanently isolated from the rest of $W$.

Therefore $\gamma(E) > 0$ is necessary for persistence under $\nu > 0$.

\medskip
\textbf{Part 2: Necessity of $\eta_M > 0$.}

Suppose $\gamma > 0$ but $\eta_M \leq 0$.

\emph{Case $\eta_M = 0$}: The entity engages $M$ but $M$ produces no
improvement in grounding. $(A, K)$ is effectively fixed despite
engagement. The free-context exhaustion argument applies: novel states
consume context that is never freed by promotion. Additionally, $C_n$ is
bounded. Promotion to lower encoding levels requires achieving
the rigor threshold $\theta_i$ (Definition~\ref{def:encoding}), which
requires $M$ to improve the relevant action to the required certainty.
With $\eta_M = 0$, no action reaches threshold, $L_n$ is never
relieved, and $\rho \to 0$.

\emph{Case $\eta_M < 0$ persistently}: $(A, K)$ worsens in expectation,
accelerating the failure mode.

The persistence condition is:
\[
\mathbb{E}_t[\eta_M] > 0 \quad \text{over the entity's trajectory}
\]
not $\eta_M(t) > 0$ at every time step. Temporary $\eta_M < 0$
(exploration, creative destruction) is compatible with persistence.

\medskip
\textbf{Part 3: Minimality.}

For every other scalar property of $E$, there exists a persistent
entity for which that property equals zero:
\begin{itemize}
\item $K(0) = \varnothing$: built from monotone descent
      (Theorem~\ref{thm:descent}).
\item $A(0) = \{\mathrm{id}\}$: the identity action still returns
      feedback from $W$.
\item $C_n = C_{\min}$: graph structure \cite{McCarthy2026a}
      maintains $\rho = 1$ by retrieving the relevant subgraph.
\end{itemize}

The exception is sensing capacity: $\gamma > 0$ requires learning
signals, which require that the selection mechanism $\sigma$ be coupled
to $W$ via sensing. But sensing is excluded from the invariants not by
definitional fiat but by structural role: it is a \emph{precondition}
for the entity to couple to $W$ at all. The tuple $(K, A, \sigma,
\gamma)$ with zero sensing reduces to a static object that neither
observes nor acts. The theorem's domain is entities that interact with
$W$, which requires positive sensing as a precondition of interaction,
not as a variable property that selection acts upon. The distinction is between
preconditions (without which the framework does not apply) and
invariants (properties that vary within the framework and whose values
are selected for).

\medskip
\textbf{Part 4: $I(E) = \gamma \cdot \eta_M$ follows by arithmetic.}

Parts~1--3 are the substantive content. Given that $\gamma$ is a
probability and $\eta_M$ a conditional rate, their product is the
expected improvement rate by the law of total expectation
(Corollary~\ref{cor:measure}).
\end{proof}

% ------------------------------------------------------------------
\section{$I(E) = \gamma \cdot \eta_M$ Is Structurally Necessary}
\label{sec:measure}
% ------------------------------------------------------------------

\begin{corollary}[The Intelligence Measure]
\label{cor:measure}
$\gamma$ and $\eta_M$ are the unique structural invariants of persistence
(Theorem~\ref{thm:invariants}). The expected improvement rate is:
\[
I(E) = \gamma \cdot \eta_M
\]
\end{corollary}

\begin{proof}
$\gamma \in [0, 1]$ is the marginal probability of engaging $M$.
$\eta_M$ is the expected improvement rate conditional on engagement.
By the law of total expectation:
\[
\mathbb{E}[\text{improvement rate}]
= \gamma \cdot \mathbb{E}[\text{rate} \mid \text{engaged}]
+ (1 - \gamma) \cdot 0
= \gamma \cdot \eta_M
\]
This identity holds whether $\gamma$ is content-blind or directed.
The conditioning in Definition~\ref{def:eta} absorbs the selection
effect. The product is exact.

The substantive claim is not the formula but the identification:
intelligence is the rate at which an entity \emph{adapts to novelty},
not the level of its current performance. The derivation makes this
precise: $\gamma > 0$ and $\eta_M > 0$ are required specifically
because $\nu(E) > 0$---novel states arrive that would be lethal
without adaptation. Intelligence is what positive novelty rate
selects for.
This distinguishes the present account from Hutter
\cite{Hutter2005} (intelligence as performance level) and from
ordinal rankings that treat intelligence as a comparison without
units. Given this identification, the product follows uniquely by
total expectation: $\gamma$ is a probability, $\eta_M$ is a
conditional rate, and their product is the unconditional rate.
Any monotone transform $f(\gamma \cdot \eta_M)$ would measure a
\emph{function of} the rate, not the rate itself, and would lack
the dimensional interpretation (bits of improvement per learning
signal). The uniqueness is not that the product is the only
function zero on the axes---many are---but that it is the only
function that \emph{is} the expected improvement rate rather than
a rescaling of it.
\end{proof}

\textbf{Temporal note.} $I(E)$ is the expected improvement per learning
signal, not per unit time. The temporal rate of improvement depends on
the entity-environment coupling: if learning signals arrive at rate
$\lambda(t)$, the instantaneous improvement rate is
$I(E) \cdot \lambda(t)$. Two entities with identical $I(E)$ in different
environments improve at different temporal rates. $I(E)$ measures the
entity's intrinsic capacity; the environment determines how often that
capacity is exercised.

\begin{remark}[Exploration and Creative Destruction]
\label{rem:exploration}
Temporarily negative $\eta_M$---locally irrational steps that reduce
grounding in the short run---is survival-optimal when the persistence
condition requires only $\mathbb{E}_t[\eta_M] > 0$ over the trajectory.
An entity restricted to $\eta_M \geq 0$ at every step is trapped in
local optima; an entity that can tolerate temporary $\eta_M < 0$
accesses a larger region of the $(A, K)$ landscape.
\end{remark}

% ------------------------------------------------------------------
\section{Engagement with Prior Work}
\label{sec:prior}
% ------------------------------------------------------------------

\subsection{Intelligence Measures and Definitions}

\textbf{Dupoux, LeCun, and Malik \cite{DupouxLeCunMalik2026}.}
They deliberately avoid defining intelligence, instead characterizing
it as the capacity for autonomous learning. Theorem~\ref{thm:invariants}
establishes that this is structurally correct: $\gamma > 0$ and
$\eta_M > 0$ are the unique invariants. Their product is the definition
they avoided giving. Their bilevel optimization framework (inner loop as
development, outer loop as evolution) captures the correct structural
relationship between encoding levels, but treats the learning rate
difference as a design parameter: the inner loop is fast because it is
designed to be fast. Definition~\ref{def:encoding} derives this from
first principles: $\eta_{\mathrm{evo}} \ll \eta_{\mathrm{ind}}$ because
$C_0 \gg C_n$. Their ``System~M'' that is ``not learned through
gradient descent'' but fixed by evolution is explained by
$\theta_0 \approx 1$: only mechanisms validated across maximal
environmental variation meet this threshold.

\textbf{Hutter \cite{Hutter2005} and Legg--Hutter \cite{LeggHutter2007}.}
Three structural differences: (1)~AIXI assumes unbounded computation;
we have bounded $C_n$. Under bounded context, the AIXI solution is
unreachable: the entity descends with finite resources
(Theorem~\ref{thm:descent}). (2)~They define intelligence as
performance level; we define it as rate of improvement, independent of
current performance. (3)~Their environment class is all computable
functions; ours is physical $W$ with positive entropy.

\textbf{Chollet \cite{Chollet2019}.}
Skill-acquisition efficiency is structurally close to $\eta_M$. Two
differences: Chollet treats prior knowledge $K$ as a confound to be
controlled for; here $K$ is constitutive, because $\eta_M$ operates
on $(A, K)$ and the efficiency of learning depends on what is already
known. Chollet's ``scope'' and ``generalization difficulty'' map to the
growth of reachable territory in the $(A, K)$ landscape.

\textbf{Schmidhuber \cite{Schmidhuber2010}.}
Curiosity as compression progress corresponds to $\eta_M$. The
difference is foundational. Schmidhuber treats intrinsic drive as a
design objective. Theorem~\ref{thm:invariants}, Part~1, establishes
$\gamma > 0$ as a necessary condition of persistence: not a reward to
be designed in, but a property that selection enforces.

\textbf{Cattell \cite{Cattell1971} and Spearman \cite{Spearman1904}.}
Cattell's fluid intelligence ($G_f$) maps to $\eta_M$; crystallized
intelligence ($G_c$) maps to $|K|$. The framework predicts the
well-established correlation: entities with higher $\eta_M$ build
larger $K$ faster. Spearman's $g$-factor should correlate with
$I(E) = \gamma \cdot \eta_M$---a testable prediction.

\subsection{Self-Organization and Minimization}

\textbf{Friston \cite{Friston2010, Friston2017}.}
The free energy principle derives minimization behavior from
self-organization: systems that persist must minimize surprise (free
energy). This is consistent with monotone descent on $H(W \mid K)$
(Theorem~\ref{thm:descent}). FEP does contain an elimination
argument: systems that fail to minimize free energy cease to exist as
identifiable systems. However, FEP does not distinguish between
$\gamma = 0$ and $\gamma > 0$ within its framework: a rock trivially
minimizes free energy, and a thermostat minimizes it with fixed
$(A, K)$. The present account adds a specific mechanism---free-context
exhaustion under positive novelty rate---that eliminates $\gamma = 0$
entities. The contribution is the context-filling argument
(Theorem~\ref{thm:invariants}, Part~1): not that minimization occurs,
but that active, learning-driven minimization ($\gamma > 0$) is
necessary when $\nu > 0$.

\subsection{Encoding Hierarchy in Practice}

\textbf{Hinton and Nowlan \cite{Hinton1987}.}
They demonstrate via simulation that individual learning can accelerate
evolutionary search: the Baldwin effect. Our framework formalizes this:
individual learning discovers actions at $L_n$; encoding permanence
(Definition~\ref{def:encoding}) determines which promote to $L_0$. The
Baldwin effect is a necessary consequence of the encoding hierarchy.

\textbf{Kirkpatrick et al.\ \cite{Kirkpatrick2017}.}
Elastic weight consolidation assigns higher importance to weights
encoding well-established knowledge, penalizing changes during
subsequent learning. This is the encoding hierarchy: weights at lower
levels (higher $\theta_i$) are protected; weights at higher levels are
free to change. EWC's Fisher information matrix empirically estimates
our $C_i$. The result provides direct evidence that artificial systems
violating the hierarchy (flat learning rate) fail, and systems
approximating it succeed.

\textbf{Murray et al.\ \cite{Murray2014}.}
Intrinsic timescales of neural activity increase along the cortical
hierarchy: sensory cortex fluctuates rapidly, prefrontal cortex
maintains stable representations. This maps onto the encoding
hierarchy: fast-updating sensory areas at $L_n$, slow-updating
prefrontal areas at lower levels. The empirical timescale gradient
matches $\eta_i \propto 1/\log(C_i / \varepsilon)$.

\textbf{Finn et al.\ \cite{FinnEtAl2017} and Thrun
\cite{Thrun1998}.}
MAML and ``learning to learn'' are $M$ operating on $M$: the
$M^{(2)}$ level. The meta-learner updates more slowly
($\eta_{M^{(2)}} < \eta_M$) because $C_{M^{(2)}} > C_M$. The MAML
inner loop is $L_n$; the outer loop is $L_{n-1}$. Evolution is
meta-meta-$\ldots$-learning at $L_0$.

% ------------------------------------------------------------------
\section{The Complete Chain}
\label{sec:chain}
% ------------------------------------------------------------------

The four papers derive the necessary architecture of intelligence from
a single generative constraint:

\begin{enumerate}
\item \textbf{Bounded $C_n$ + information preservation $\;\Rightarrow\;$
      graph structure} \cite{McCarthy2026a}. Any entity accumulating
      propositions under bounded context necessarily maintains a
      directed graph with typed edges.

\item \textbf{Bounded $C_n$ + growing $A$ $\;\Rightarrow\;$ convergence
      to escalation} \cite{McCarthy2026b}. Any viable control
      architecture converges toward escalation properties as concurrent
      demands grow relative to $C_n$.

\item \textbf{Bounded context $\;\Rightarrow\;$ normalization of shared
      propositions into $K$, with asymmetric retention}
      \cite{McCarthy2026c}. Overlapping action policies under bounded
      $C_n$ force extraction of shared propositions into separated $K$;
      the retention criterion is probability of need $\times$ conditional
      value.

\item \textbf{$\gamma > 0$ + feedback $\;\Rightarrow\;$ descent on $U$
      + encoding hierarchy; positive entropy $\;\Rightarrow\;$
      $\gamma > 0$ and $\eta_M > 0$ necessary;
      $I(E) = \gamma \cdot \eta_M$} [this paper]. Agency with feedback
      produces convergent uncertainty reduction. Cost-of-failure
      generates the encoding hierarchy. The lineage and context-filling
      arguments establish the two invariants; their product is the unique
      intelligence measure.
\end{enumerate}

Each step is a necessary consequence of bounded context interacting
with properties of the physical world: positive entropy, survival
pressure, and feedback.

% ------------------------------------------------------------------
\section{Open Questions}
\label{sec:open}
% ------------------------------------------------------------------

\begin{enumerate}
\item \textbf{Origin of $\gamma$.} Theorem~\ref{thm:invariants}
      establishes that $\gamma > 0$ is necessary but does not explain
      how $\gamma > 0$ structures bootstrap from $\gamma = 0$
      substrates. The thermodynamic conjecture: dissipative structures
      \cite{Prigogine1984, England2013} under sustained energy flow
      spontaneously produce $\gamma > 0$ configurations, which selection
      retains.

\item \textbf{Direction of $\gamma$.} Is agency uniformly distributed
      across learning signals, or does it have intrinsic directedness?

\item \textbf{The $M^{(n)}$ hierarchy.} How many levels of
      meta-learning does a persistent entity require? Is there a natural
      termination point, or is the hierarchy open upward without bound?

\item \textbf{Bound composition in $A$.} How do individual action
      bounds compose when multiple actions share knowledge substructure?
\end{enumerate}

% ------------------------------------------------------------------
\section{Falsifiable Predictions}
\label{sec:predictions}
% ------------------------------------------------------------------

The theorems yield eight predictions that are empirically testable.

\textbf{1. Catastrophic forgetting will not be solved within flat
parameter spaces.}
Theorem~\ref{thm:type} predicts that any continual learning solution
that works at scale will implement $A$/$M$ stratification---whether
explicitly (elastic weight consolidation, progressive networks,
complementary learning systems) or implicitly (by learning to protect
certain parameters from learning updates). A solution that uses a
single undifferentiated parameter space for both action execution and
learning, with no form of weight protection or functional separation,
will fail on long-horizon continual learning benchmarks (100+
sequential tasks). A flat-parameter system that matches stratified
architectures on such benchmarks without any form of functional
separation would falsify Theorem~\ref{thm:type}.

\textbf{2. Effective agent architectures will converge to escalation.}
Multi-step AI agents (planning-based systems that decompose goals into
subtask sequences) operate under bounded $C_n$: each delegation
requires tracking context. Theorem~\ref{thm:invariants} and the
encoding hierarchy (Definition~\ref{def:encoding}) predict that
successful agent architectures will converge toward escalation:
routine subtasks handled by compiled, cached, or tool-based policies
at lower encoding levels, with the frontier model invoked only on
failure or novelty. Purely top-down planning agents that maintain
full supervisory state over every subtask will fail to scale beyond
$k^* = C_n / (2\alpha)$ concurrent delegations. A purely top-down
planning agent that scales to 1000+ step tasks without delegating to
compiled subroutines would falsify the escalation convergence result.

\textbf{3. $I(E) = \gamma \cdot \eta_M$ predicts $G_f$ correlations.}
If intelligence is the rate of improvement, then Cattell's fluid
intelligence $G_f$ (which measures novel problem-solving capacity)
should correlate more strongly with $\eta_M$ than with $|K|$. The
prediction: across individuals, $G_f$ scores will predict learning
rate on novel tasks better than they predict current knowledge
volume. Crystallized intelligence $G_c$ should correlate with
$|K|$ and with the cumulative integral of $I(E)$ over time. A
population in which $G_c$ predicts novel-task learning rate better
than $G_f$ does would falsify the identification.

\textbf{4. $\gamma = 0$ systems fail under distribution shift.}
The free-context exhaustion argument
(Theorem~\ref{thm:invariants}, Part~1) predicts that systems with
$\gamma = 0$ (no capacity to update $K$ or $A$ in response to
feedback) will fail under sustained distribution shift, even if they
perform well on the initial distribution. The mechanism is specific:
$C_{\mathrm{free}} \to 0$ as unresolvable novel states accumulate.
The prediction: deployed ML systems without online learning
capabilities will show monotonically degrading performance under
distribution drift, with the degradation rate proportional to the
novelty rate $\nu$. A static model that maintains performance
indefinitely under sustained distribution shift would falsify the
free-context exhaustion argument.

\textbf{5. Graph-structured knowledge will produce cross-specialty
unifications.}
The graph necessity result of \cite{McCarthy2026a}
establishes that shared propositions are factored into reusable nodes
with directed typed edges. When knowledge from different specialties
is encoded in a shared graph, propositions that were independently
maintained in separate domains become structurally adjacent---their
shared dependencies become visible. The prediction: as institutional
and scientific knowledge graphs grow, cross-specialty unifications
(where two fields discover they are studying the same underlying
structure under different names) will increase in frequency, driven
not by individual insight but by the graph making shared structure
visible. Falsified if large-scale knowledge graphs produce no
increase in detected cross-domain structural equivalences compared
to unstructured literature search.

\textbf{6. Open knowledge graphs with provenance will undercut
parametric knowledge storage.}
The normalization result \cite{McCarthy2026c} establishes that
commonly held knowledge---the high-overlap region of $K$ across
entities---is exactly the set of propositions that factoring extracts
into shared, reusable nodes. Provenance is a typed edge:
each proposition carries its reconstruction path to source evidence
\cite{McCarthy2026a}. An open-source knowledge graph encoding commonly
held propositions with provenance would provide $O(d)$ retrieval with
full attribution at near-zero marginal cost. Current large language
models store the same commonly held knowledge in continuous weights:
retrieval is $O(n^2)$, provenance is absent, and the knowledge is
neither inspectable nor correctable. The prediction: open knowledge
graphs with provenance will replicate and eventually surpass the
factual-knowledge function of parametric models, at a fraction of
the compute cost, because structured retrieval with attribution is
strictly cheaper than re-deriving answers from weights on every query.
Falsified if parametric storage achieves cost parity with structured
retrieval on factual knowledge tasks at arbitrary scale while
simultaneously providing equivalent provenance.

\textbf{7. The first widely recognized artificial general intelligence
will have explicit graph structure, not monolithic parameters.}
Theorem~\ref{thm:invariants} identifies the structural invariants of
intelligence as $\gamma > 0$ (agency) and $\eta_M > 0$ (learning
efficiency). Current large language models have $\gamma = 0$: they do
not update $K$ or $A$ in response to feedback during deployment.
They approximate intelligence by storing a large static $K$ in weights
and simulating $\gamma > 0$ within the context window---a strategy
the free-context exhaustion argument
(Theorem~\ref{thm:invariants}, Part~1) predicts cannot persist under
positive novelty rate. The framework predicts that the first system
widely recognized as artificial general intelligence will exhibit:
(i)~explicit or learned graph structure over propositions
\cite{McCarthy2026a}, (ii)~escalation-based control
\cite{McCarthy2026b} rather than monolithic inference,
(iii)~separated $K$ with provenance rather than knowledge fused into
parameters, and (iv)~$\gamma > 0$---genuine online learning that
updates $(A, K)$ from feedback. These are not design recommendations
but convergence results: any system that persists under bounded context
and positive novelty rate will exhibit these properties regardless of
substrate. Falsified if a system with monolithic parameters, no graph
structure, no knowledge separation, and no online learning achieves
robust, open-ended performance across novel domains over extended
deployment without degradation.

\textbf{8. Cancer treatment will converge on restoring broken
escalation chains.}
The escalation hierarchy \cite{McCarthy2026b} predicts that
multicellular organisms handle errors at the lowest sufficient level:
DNA repair ($L_0$), cell-cycle arrest ($L_1$), apoptosis ($L_2$),
immune surveillance ($L_3$), with conscious medical intervention at
$L_n$. Cancer is the pathology of broken escalation: mutations disable
one or more lower-level error handlers, so cellular errors that should
be caught and resolved locally instead escalate unchecked to the
organism level. Current treatment paradigms attack the cancer
directly at $L_n$ (surgery, chemotherapy, radiation)---the equivalent
of top-down control. The escalation framework predicts that effective
treatment will converge toward restoring the broken levels of the
chain rather than compensating at $L_n$. Checkpoint immunotherapy
(restoring immune surveillance), synthetic lethality approaches
(exploiting remaining repair pathways), and p53 reactivation therapies
are early instances: each restores a lower-level handler rather than
attacking the tumor directly. The prediction: treatment efficacy will
correlate with the number of escalation levels restored, not with the
directness of tumor targeting, and the field will shift from
``kill the cancer'' toward ``repair the chain that should have caught
it.'' Falsified if direct tumor-targeting therapies consistently
outperform escalation-restoring therapies across cancer types in
long-term survival outcomes.

% ------------------------------------------------------------------
\section{Conclusion}
\label{sec:conclusion}
% ------------------------------------------------------------------

Intelligence is the rate of adaptation to novelty. This is not proposed
but identified by elimination: positive novelty rate requires $\gamma > 0$
and $\eta_M > 0$; their product is the unconditional rate.

Theorem~\ref{thm:descent} establishes that agency with feedback
produces monotone local descent on expected uncertainty $H(W \mid K)$.
Each informative update reduces expected uncertainty by at most the
information gain, with equality under exact Bayesian conditioning. Under
non-stationary task distribution, descent is local and instantaneous
rather than globally convergent.

Definition~\ref{def:encoding} derives the encoding hierarchy from
cost-of-failure: $\theta_i = 1 - \varepsilon / C_i$ determines the
certainty required for promotion, and $\eta_{\mathrm{evo}} \ll
\eta_{\mathrm{ind}}$ follows from $C_0 \gg C_n$.

Theorem~\ref{thm:type} establishes that $A$/$M$ stratification is
selectively favored under bounded context: collapsing action execution
and learning into a single parameter space is possible but increasingly
costly under positive novelty rate.

Theorem~\ref{thm:invariants} establishes that $\gamma > 0$ and
$\eta_M > 0$ are the unique structural invariants of persistence under
positive novelty rate. The free-context exhaustion argument eliminates
$\gamma = 0$ entities; the promotion argument eliminates $\eta_M = 0$
entities. The product $I(E) = \gamma \cdot \eta_M$ is the unique measure
consistent with these invariants, derived from the law of total
expectation applied to a probability gate and a conditional rate.

The four-paper chain derives, from the single constraint of bounded
active context, the necessary architecture: graph structure,
escalation, knowledge--action inseparability, evaluation-driven descent
with encoding permanence, and the structural invariants of intelligence.
Each result is a necessary consequence of bounded context interacting
with properties of the physical world: positive entropy, survival
pressure, and feedback.

\section*{Acknowledgments}

The author used Claude (Anthropic) for drafting assistance, literature review, and \LaTeX{} preparation. All theoretical content, proofs, and arguments are the author's own.

\bibliographystyle{plain}
\begin{thebibliography}{99}

\bibitem{Ashby1956}
W.~R. Ashby.
\newblock \emph{An Introduction to Cybernetics}.
\newblock Chapman and Hall, 1956.

\bibitem{Barendregt1984}
H.~P. Barendregt.
\newblock \emph{The Lambda Calculus: Its Syntax and Semantics}.
\newblock North-Holland, revised edition, 1984.

\bibitem{Cattell1971}
R.~B. Cattell.
\newblock \emph{Abilities: Their Structure, Growth, and Action}.
\newblock Houghton Mifflin, 1971.

\bibitem{Chollet2019}
F.~Chollet.
\newblock On the measure of intelligence.
\newblock \emph{arXiv preprint arXiv:1911.01547}, 2019.

\bibitem{Cover2006}
T.~M. Cover and J.~A. Thomas.
\newblock \emph{Elements of Information Theory}.
\newblock Wiley-Interscience, 2nd edition, 2006.

\bibitem{DupouxLeCunMalik2026}
E.~Dupoux, Y.~LeCun, and J.~Malik.
\newblock Why {AI} systems don't learn and what to do about it: Lessons on
  autonomous learning from cognitive science.
\newblock \emph{arXiv preprint arXiv:2603.15381}, 2026.

\bibitem{England2013}
J.~L. England.
\newblock Statistical physics of self-replication.
\newblock \emph{The Journal of Chemical Physics}, 139(12):121923, 2013.

\bibitem{FinnEtAl2017}
C.~Finn, P.~Abbeel, and S.~Levine.
\newblock Model-agnostic meta-learning for fast adaptation of deep networks.
\newblock In \emph{Proceedings of the 34th International Conference on Machine
  Learning}, pages 1126--1135, 2017.

\bibitem{Friston2010}
K.~Friston.
\newblock The free-energy principle: A unified brain theory?
\newblock \emph{Nature Reviews Neuroscience}, 11(2):127--138, 2010.

\bibitem{Friston2017}
K.~Friston.
\newblock Active inference and learning.
\newblock \emph{Neuroscience and Biobehavioral Reviews}, 68:862--879, 2017.

\bibitem{Hinton1987}
G.~E. Hinton and S.~J. Nowlan.
\newblock How learning can guide evolution.
\newblock \emph{Complex Systems}, 1(3):495--502, 1987.

\bibitem{Hoeffding1963}
W.~Hoeffding.
\newblock Probability inequalities for sums of bounded random variables.
\newblock \emph{Journal of the American Statistical Association},
  58(301):13--30, 1963.

\bibitem{Hutter2005}
M.~Hutter.
\newblock \emph{Universal Artificial Intelligence: Sequential Decisions Based
  on Algorithmic Probability}.
\newblock Springer, 2005.

\bibitem{Kirkpatrick2017}
J.~Kirkpatrick, R.~Pascanu, N.~Rabinowitz, J.~Veness, G.~Desjardins,
  A.~A. Rusu, K.~Milan, J.~Quan, T.~Ramalho, A.~Grabska-Barwinska,
  D.~Hassabis, C.~Summerfield, T.~Lillicrap, and D.~Kumaran.
\newblock Overcoming catastrophic forgetting in neural networks.
\newblock \emph{Proceedings of the National Academy of Sciences},
  114(13):3521--3526, 2017.

\bibitem{McClelland1995}
J.~L. McClelland, B.~L. McNaughton, and R.~C. O'Reilly.
\newblock Why there are complementary learning systems in the hippocampus
  and neocortex: Insights from the successes and failures of connectionist
  models of learning and memory.
\newblock \emph{Psychological Review}, 102(3):419--457, 1995.

\bibitem{LeggHutter2007}
S.~Legg and M.~Hutter.
\newblock Universal intelligence: A definition of machine intelligence.
\newblock \emph{Minds and Machines}, 17(4):391--444, 2007.

\bibitem{McCarthy2026a}
P.~D. McCarthy.
\newblock Graph structure is necessary for information preservation under
  bounded context.
\newblock Manuscript, 2026.

\bibitem{McCarthy2026b}
P.~D. McCarthy.
\newblock Convergence of control architectures to escalation under bounded
  context.
\newblock Manuscript, 2026.

\bibitem{McCarthy2026c}
P.~D. McCarthy.
\newblock Knowledge as normalization: bounded context separates shared
  propositions from action.
\newblock Manuscript, 2026.

\bibitem{Murray2014}
J.~D. Murray, A.~Bernacchia, D.~J. Freedman, R.~Romo, J.~D. Wallis,
  X.~Cai, C.~Padoa-Schioppa, T.~Pasternak, H.~Seo, D.~Lee, and
  X.-J. Wang.
\newblock A hierarchy of intrinsic timescales across primate cortex.
\newblock \emph{Nature Neuroscience}, 17(12):1661--1663, 2014.

\bibitem{Prigogine1984}
I.~Prigogine and I.~Stengers.
\newblock \emph{Order Out of Chaos: Man's New Dialogue with Nature}.
\newblock Bantam Books, 1984.

\bibitem{Schmidhuber2010}
J.~Schmidhuber.
\newblock Formal theory of creativity, fun, and intrinsic motivation
  (1990--2010).
\newblock \emph{IEEE Transactions on Autonomous Mental Development},
  2(3):230--247, 2010.

\bibitem{Shannon1948}
C.~E. Shannon.
\newblock A mathematical theory of communication.
\newblock \emph{Bell System Technical Journal}, 27(3):379--423, 1948.

\bibitem{Spearman1904}
C.~Spearman.
\newblock ``General intelligence,'' objectively determined and measured.
\newblock \emph{The American Journal of Psychology}, 15(2):201--292, 1904.

\bibitem{Thrun1998}
S.~Thrun and T.~M. Mitchell.
\newblock Learning to learn.
\newblock In \emph{Learning to Learn}, pages 3--17. Springer, 1998.

\end{thebibliography}

\end{document}
